% GENERATED FILE. Edit canonical lessons, then run npm run generate:books.
% canonical-source-hash: fnv1a64:c3c94d6d4935e3fc
% canonical-lessons: HI-C275-moc, HI-C275-hicki, HI-C275-jhapki, HI-C275-karvat, HI-C275-angrai, HI-R275-first-pass-words-from-chapters-270-272, HI-R275-second-pass-words-from-chapters-273-275

\chapter{A sprain, A hiccup, A nap, Turning over, A stretch}
\label{ch:hi-a-sprain-a-hiccup-a-nap-turning-over-a-stretch}
\hlchaptermodality{275}

\begin{chapteropening}
\textbf{By the end of this chapter:} \emph{I can say a sprain, a hiccup, a nap, turning over and a stretch.}

Complete the last lesson of chapter 275: I can say a sprain, a hiccup, a nap, turning over and a stretch.
\end{chapteropening}

\section[\texorpdfstring{\hi{मोच} (moc)}{मोच (moc)}]{\hi{मोच} (moc) — a sprain}

\label{lesson:HI-C275-moc}



\begin{quote}
Before the new one: say the Hindi for a burning feeling, then the Hindi for a swelling.
\end{quote}

\subsection*{You'll want to know: \hi{मोच}}
\textbf{\hi{मोच}} — \emph{moc} — \textquotedblleft{}a sprain\textquotedblright{}.

\begin{grammarlens}[title={What you've built}]
One more word for this chapter.
\end{grammarlens}

\subsection*{Your turn}
\begin{itemize}
\raggedright
  \item \emph{Say it:} \emph{moc}
  \item \emph{Say it:} \emph{moc}, once more
  \item \emph{Say it:} say \emph{moc}
  \item \emph{Recall:} say the Hindi for a burning feeling, then the Hindi for a swelling, then say \emph{moc} again
\end{itemize}

\begin{tcolorbox}[breakable,colback=teal!4,colframe=teal!35!black,title={Before you move on}]
What does \hi{मोच} mean? (\textquotedblleft{}A sprain\textquotedblright{}.) Say it once more.
\end{tcolorbox}
\section[\texorpdfstring{\hi{हिचकी} (hickī)}{हिचकी (hickī)}]{\hi{हिचकी} (hickī) — a hiccup}

\label{lesson:HI-C275-hicki}



\begin{quote}
Before the new one: say the Hindi for a swelling, then the Hindi for a sprain.
\end{quote}

\subsection*{You'll want to know: \hi{हिचकी}}
\textbf{\hi{हिचकी}} — \emph{hickī} — \textquotedblleft{}a hiccup\textquotedblright{}.

\begin{grammarlens}[title={What you've built}]
One more word for this chapter.
\end{grammarlens}

\subsection*{Your turn}
\begin{itemize}
\raggedright
  \item \emph{Say it:} \emph{hickī}
  \item \emph{Say it:} \emph{hickī}, once more
  \item \emph{Say it:} say \emph{hickī}
  \item \emph{Recall:} say the Hindi for a swelling, then the Hindi for a sprain, then say \emph{hickī} again
\end{itemize}

\begin{tcolorbox}[breakable,colback=teal!4,colframe=teal!35!black,title={Before you move on}]
What does \hi{हिचकी} mean? (\textquotedblleft{}A hiccup\textquotedblright{}.) Say it once more.
\end{tcolorbox}
\section[\texorpdfstring{\hi{झपकी} (jhapkī)}{झपकी (jhapkī)}]{\hi{झपकी} (jhapkī) — a nap}

\label{lesson:HI-C275-jhapki}



\begin{quote}
Before the new one: say the Hindi for a sprain, then the Hindi for a hiccup.
\end{quote}

\subsection*{You'll want to know: \hi{झपकी}}
\textbf{\hi{झपकी}} — \emph{jhapkī} — \textquotedblleft{}a nap\textquotedblright{}.

\begin{grammarlens}[title={What you've built}]
One more word for this chapter.
\end{grammarlens}

\subsection*{Your turn}
\begin{itemize}
\raggedright
  \item \emph{Say it:} \emph{jhapkī}
  \item \emph{Say it:} \emph{jhapkī}, once more
  \item \emph{Say it:} say \emph{jhapkī}
  \item \emph{Recall:} say the Hindi for a sprain, then the Hindi for a hiccup, then say \emph{jhapkī} again
\end{itemize}

\begin{tcolorbox}[breakable,colback=teal!4,colframe=teal!35!black,title={Before you move on}]
What does \hi{झपकी} mean? (\textquotedblleft{}A nap\textquotedblright{}.) Say it once more.
\end{tcolorbox}
\section[\texorpdfstring{\hi{करवट} (karvaṭ)}{करवट (karvaṭ)}]{\hi{करवट} (karvaṭ) — turning over (in bed)}

\label{lesson:HI-C275-karvat}



\begin{quote}
Before the new one: say the Hindi for a hiccup, then the Hindi for a nap.
\end{quote}

\subsection*{You'll want to know: \hi{करवट}}
\textbf{\hi{करवट}} — \emph{karvaṭ} — \textquotedblleft{}turning over (in bed)\textquotedblright{}.

\begin{grammarlens}[title={What you've built}]
One more word for this chapter.
\end{grammarlens}

\subsection*{Your turn}
\begin{itemize}
\raggedright
  \item \emph{Say it:} \emph{karvaṭ}
  \item \emph{Say it:} \emph{karvaṭ}, once more
  \item \emph{Say it:} say \emph{karvaṭ}
  \item \emph{Recall:} say the Hindi for a hiccup, then the Hindi for a nap, then say \emph{karvaṭ} again
\end{itemize}

\begin{tcolorbox}[breakable,colback=teal!4,colframe=teal!35!black,title={Before you move on}]
What does \hi{करवट} mean? (\textquotedblleft{}Turning over (in bed)\textquotedblright{}.) Say it once more.
\end{tcolorbox}
\section[\texorpdfstring{\hi{अंगड़ाई} (aṅgṛāī)}{अंगड़ाई (aṅgṛāī)}]{\hi{अंगड़ाई} (aṅgṛāī) — a stretch (of the body)}

\label{lesson:HI-C275-angrai}



\begin{quote}
Before the new one: say the Hindi for a nap, then the Hindi for turning over.
\end{quote}

\subsection*{You'll want to know: \hi{अंगड़ाई}}
\textbf{\hi{अंगड़ाई}} — \emph{aṅgṛāī} — \textquotedblleft{}a stretch (of the body)\textquotedblright{}.

\begin{grammarlens}[title={What you've built}]
One more word for this chapter.
\end{grammarlens}

\subsection*{Your turn}
\begin{itemize}
\raggedright
  \item \emph{Say it:} \emph{aṅgṛāī}
  \item \emph{Say it:} \emph{aṅgṛāī}, once more
  \item \emph{Say it:} say \emph{aṅgṛāī}
  \item \emph{Recall:} say the Hindi for a nap, then the Hindi for turning over, then say \emph{aṅgṛāī} again
\end{itemize}

\begin{tcolorbox}[breakable,colback=teal!4,colframe=teal!35!black,title={Before you move on}]
What does \hi{अंगड़ाई} mean? (\textquotedblleft{}A stretch (of the body)\textquotedblright{}.) Say it once more.
\end{tcolorbox}
\section[(dialogue)]{Words from chapters 270-272 — a first pass}

\label{lesson:HI-R275-first-pass-words-from-chapters-270-272}



\begin{quote}
Say the words for turning over and a stretch.
\end{quote}

\subsection*{Your turn}
Say these aloud:

\begin{itemize}
\raggedright
  \item a grater, a lid, a jug, a water jar, a clay pot
  \item an office, insurance, interest, a tenant, a neighbourhood
  \item a lane, a street corner, a fair, a procession, a speech
\end{itemize}

\begin{tcolorbox}[breakable,colback=teal!4,colframe=teal!35!black,title={Before you move on}]
A grater? (\textbf{\hi{कद्दूकस}}, \emph{kaddūkas}.) A lid? (\textbf{\hi{ढक्कन}}, \emph{dhakkan}.) A jug? (\textbf{\hi{जग}}, \emph{jag}.) A water jar? (\textbf{\hi{सुराही}}, \emph{surāhī}.) A clay pot? (\textbf{\hi{मटका}}, \emph{maṭkā}.) An office? (\textbf{\hi{कार्यालय}}, \emph{kāryālay}.) Insurance? (\textbf{\hi{बीमा}}, \emph{bīmā}.) Interest? (\textbf{\hi{ब्याज}}, \emph{byāj}.) A tenant? (\textbf{\hi{किरायेदार}}, \emph{kirāyedār}.) A neighbourhood? (\textbf{\hi{मोहल्ला}}, \emph{mohallā}.) A lane? (\textbf{\hi{गली}}, \emph{galī}.) A street corner? (\textbf{\hi{नुक्कड़}}, \emph{nukkaṛ}.) A fair? (\textbf{\hi{मेला}}, \emph{melā}.) A procession? (\textbf{\hi{जुलूस}}, \emph{jalūs}.) A speech? (\textbf{\hi{भाषण}}, \emph{bhāṣaṇ}.)
\end{tcolorbox}
\section[(dialogue)]{Words from chapters 273-275 — a second pass}

\label{lesson:HI-R275-second-pass-words-from-chapters-273-275}



\begin{quote}
Say the words for the police and a pharmacy.
\end{quote}

\subsection*{Your turn}
Say these aloud:

\begin{itemize}
\raggedright
  \item the police, a pharmacy, a patient, a disease, an illness
  \item a heartbeat, breath, an itch, a burning feeling, a swelling
  \item a sprain, a hiccup, a nap, turning over, a stretch
\end{itemize}

\begin{tcolorbox}[breakable,colback=teal!4,colframe=teal!35!black,title={Before you move on}]
The police? (\textbf{\hi{पुलिस}}, \emph{pulis}.) A pharmacy? (\textbf{\hi{दवाख़ाना}}, \emph{davākhānā}.) A patient? (\textbf{\hi{रोगी}}, \emph{rogī}.) A disease? (\textbf{\hi{रोग}}, \emph{rog}.) An illness? (\textbf{\hi{बीमारी}}, \emph{bīmārī}.) A heartbeat? (\textbf{\hi{धड़कन}}, \emph{dhaṛkan}.) Breath? (\textbf{\hi{साँस}}, \emph{sā̃s}.) An itch? (\textbf{\hi{खुजली}}, \emph{khujlī}.) A burning feeling? (\textbf{\hi{जलन}}, \emph{jalan}.) A swelling? (\textbf{\hi{सूजन}}, \emph{sūjan}.) A sprain? (\textbf{\hi{मोच}}, \emph{moc}.) A hiccup? (\textbf{\hi{हिचकी}}, \emph{hickī}.) A nap? (\textbf{\hi{झपकी}}, \emph{jhapkī}.) Turning over? (\textbf{\hi{करवट}}, \emph{karvaṭ}.) A stretch? (\textbf{\hi{अंगड़ाई}}, \emph{aṅgṛāī}.)
\end{tcolorbox}
